\section{Output-Space Geometry}
\label{sec:output-geometry}

The scalar CLOC construction and its OPG matrix require no output-space metric. The level-set construction does: for a symmetric positive-definite primal metric $G_n$, output displacements have inner product $\langle u,v\rangle_{G_n}=u^\top G_n v$, the metric-normal direction is $G_n^{-1}r_n$, and the normal-distance denominator is $r_n^\top G_n^{-1}r_n$. A valid geometry therefore supplies the matched pair $G_n$ on primal output displacements and $G_n^{-1}$ on loss covectors.

The canonical choice $G_n=I$ recovers iEF \citep{IEF}. More generally, any smooth output discrepancy whose second-order expansion supplies a valid local quadratic metric can enter IEWC through its matched primal/dual pair $G_n,G_n^{-1}$. For singular ambient metrics, this requires a stated tangent space or quotient, Moore--Penrose pseudoinverse, and range conditions; diagonalizing the parameter pullback can further discard cross-parameter geometry.

For ordered probability outputs, Appendix~\ref{app:ordered-cdf-geometry} constructs a CDF/Cram\'er metric and its matched pseudoinverse dual on the simplex tangent. This supplies the space-dependent geometry that motivated the earlier diffusion experiment. That experiment measures the effect of output representation, while the appendix provides the mathematically matched geometric example.
